\section{Monte Carlo Design}

Chapter~3 uses $N$ as generic sample-size notation; the simulation chapters use $n$ for the design sample size.

\subsection{Data-Generating Process and Instrument-Relevance Extension}

The Monte Carlo experiment builds on the high-dimensional simulation design of \textcite{chen_huang_tien_2021}, with $p=100$ observed controls. For each observation $i$, the structural disturbances satisfy
\begin{equation}
  \begin{pmatrix}
    u_i \\
    \epsilon_i
  \end{pmatrix}
  \sim
  \mathcal{N}\!\left(
    \begin{pmatrix} 0 \\ 0 \end{pmatrix},
    \begin{pmatrix} 1 & 0.3 \\ 0.3 & 1 \end{pmatrix}
  \right).
  \label{eq:dgp-disturbances}
\end{equation}
The disturbance $\epsilon_i$ enters the treatment equation, whereas $u_i$ enters the outcome through $D_i u_i$. Their positive correlation therefore generates endogeneity.

The remaining primitives are mutually independent and independent of the disturbance pair in Equation~\eqref{eq:dgp-disturbances}, and they satisfy
\begin{equation}
  x_{ji}\overset{\mathrm{ind}}{\sim}\mathcal{N}(0,1),
  \quad j=1,\ldots,100,
  \qquad
  z_{1i},z_{2i},v_{1i},v_{2i}
  \overset{\mathrm{ind}}{\sim}\mathcal{N}(0,1).
  \label{eq:dgp-primitives}
\end{equation}
Observations are independent across $i$. The observed controls are
\begin{equation}
  X_{ji}=\Phi(x_{ji}),
  \qquad j=1,\ldots,100,
  \label{eq:dgp-controls}
\end{equation}
where $\Phi$ is the standard-normal cumulative distribution function, so $X_{ji}\sim\mathcal{U}(0,1)$ independently across controls.

The executable design constructs two instruments as\footnote{The displayed Monte Carlo equations in \textcite{chen_huang_tien_2021} write the corresponding control terms as $x_j$. The released simulation code instead uses the transformed controls $X_j=\Phi(x_j)$. The simulations in this thesis follow that executable implementation.}
\begin{equation}
  Z_{1i}=z_{1i}+v_{1i}+X_{2i}+X_{3i}+X_{4i},
  \label{eq:dgp-instrument1}
\end{equation}
\begin{equation}
  Z_{2i}=z_{2i}+v_{2i}+X_{7i}+X_{8i}+X_{9i}+X_{10i}.
  \label{eq:dgp-instrument2}
\end{equation}
The $z$ components enter the treatment index directly, the included controls induce dependence between instruments and controls, and the $v$ components add independent noise. The executable baseline treatment is
\begin{equation}
  D_i=\Phi(z_{1i}+z_{2i}+\epsilon_i).
  \label{eq:dgp-treatment-baseline}
\end{equation}
The treatment is continuous and lies in $(0,1)$.

The structural outcome is generated according to
\begin{equation}
  Y_i=1+D_i+5\sum_{j=1}^{7}X_{ji}+D_i u_i.
  \label{eq:dgp-outcome}
\end{equation}
The outcome contains $X_{1i},\ldots,X_{7i}$; together with the controls in the instruments, $X_{1i},\ldots,X_{10i}$ form the ten-control oracle set. For $d\in(0,1)$, Equation~\eqref{eq:dgp-outcome} implies
\begin{equation}
  Y_i(d)=1+5\sum_{j=1}^{7}X_{ji}+d(1+u_i).
  \label{eq:dgp-potential-outcome}
\end{equation}
Because $u_i$ is independent of $X_i$ and standard normal,
\begin{equation}
  Q_{Y_i(d)\mid X_i}(\tau)
  =1+5\sum_{j=1}^{7}X_{ji}
  +d\left[1+\Phi^{-1}(\tau)\right].
  \label{eq:dgp-structural-quantile}
\end{equation}
The structural quantile treatment effect is therefore
\begin{equation}
  \alpha_0(\tau)=1+\Phi^{-1}(\tau).
  \label{eq:dgp-alpha-true}
\end{equation}
In particular, $\alpha_0(0.50)=1$.

The extension varies excluded-instrument relevance through $\kappa\in[0,1]$. Let $w_i\sim\mathcal{N}(0,1)$ be independent of all preceding primitives. The latent index and treatment are
\begin{equation}
  d_i^*(\kappa)
  =\kappa(z_{1i}+z_{2i})
  +\sqrt{2(1-\kappa^2)}\,w_i
  +\epsilon_i,
  \label{eq:dgp-kappa-index}
\end{equation}
\begin{equation}
  D_i(\kappa)=\Phi\!\left(d_i^*(\kappa)\right).
  \label{eq:dgp-treatment-kappa}
\end{equation}
This symmetric loading is the sole structural modification to the executable baseline.

At $\kappa=1$, the coefficient on $w_i$ is zero and
\begin{equation}
  D_i(1)=\Phi(z_{1i}+z_{2i}+\epsilon_i),
  \label{eq:dgp-kappa-baseline}
\end{equation}
which reproduces Equation~\eqref{eq:dgp-treatment-baseline}. Independence and unit variances imply
\begin{equation}
  \operatorname{Var}\!\left[d_i^*(\kappa)\right]
  =2\kappa^2+2(1-\kappa^2)+1=3.
  \label{eq:dgp-kappa-variance}
\end{equation}
Thus $d_i^*(\kappa)\sim\mathcal{N}(0,3)$ for every $\kappa$, preserving the marginal distribution of $D_i(\kappa)$. This distribution is not uniform because $\Phi$ is applied to an index with variance three.

Instrument relevance changes because
\begin{equation}
  \operatorname{Cov}\!\left(Z_{1i},d_i^*(\kappa)\right)
  =\operatorname{Cov}\!\left(Z_{2i},d_i^*(\kappa)\right)
  =\kappa.
  \label{eq:dgp-kappa-relevance}
\end{equation}
All other instrument components are independent of $d_i^*(\kappa)$, so lower $\kappa$ reduces each instrument--index covariance. The parameter governs relevance but is not a complete measure of IVQR identification. Endogeneity remains fixed because $u_i$ is independent of $z_{1i}$, $z_{2i}$, and $w_i$ while its covariance with $\epsilon_i$ is unchanged:
\begin{equation}
  \operatorname{Cov}\!\left(u_i,d_i^*(\kappa)\right)
  =\operatorname{Cov}(u_i,\epsilon_i)=0.3.
  \label{eq:dgp-kappa-endogeneity}
\end{equation}
Joint normality and the invariant variance and covariance in Equations~\eqref{eq:dgp-kappa-variance} and~\eqref{eq:dgp-kappa-endogeneity} keep the joint distribution of $(u_i,d_i^*(\kappa))$ unchanged. The controls, observed-instrument construction, structural outcome equation, induced marginal outcome distribution, and structural quantile effect are also unchanged. The extension therefore isolates relevance from the marginal treatment distribution, endogeneity channel, and structural treatment-effect distribution.

\subsection{Simulation Design and Estimator Comparison}

The completed Monte Carlo experiment uses the design dimensions
\begin{equation}
  \begin{aligned}
    n&\in\{500,1000\},
    &\tau&\in\{0.10,0.25,0.50,0.75,0.90\},\\
    \kappa&\in\{1,0.50,0.25,0.10\},
    &R&=500.
  \end{aligned}
  \label{eq:simulation-design}
\end{equation}
The Cartesian product contains $40$ distinct $(n,\tau,\kappa)$ designs, each with $500$ replications.

Oracle-GMM uses the known oracle control set; Full-GMM uses all controls without regularization; and DML-IVQR uses regularized high-dimensional nuisances and an orthogonalization-based score construction, with the implementation-specific residualization intercept described in Section~3.3. The design therefore has $120$ design--estimator cells. Every comparison uses the truth $\alpha_0(\tau)=1+\Phi^{-1}(\tau)$ from Equation~\eqref{eq:dgp-alpha-true}.

All estimators use the same sample within a replication. DML-IVQR uses the BC pivotal penalty and same-sample procedure in Section~3.3. Inference is conducted at 95\%, and power uses the alternatives defined below.

\subsection{Performance Measures and Numerical Inference}

For a fixed $(n,\tau,\kappa,\text{estimator})$ cell, let $\widehat{\alpha}_r(\tau)$ denote the point estimate in replication $r$, and let $\alpha_0(\tau)$ be the true value from Equation~\eqref{eq:dgp-alpha-true}. Replications are indexed by $r=1,\ldots,R$. Suppressing an estimator superscript, point-estimation performance is summarized by
\begin{equation}
  \begin{aligned}
    \operatorname{Bias}(\tau)
      &=\frac{1}{R}\sum_{r=1}^{R}\left[\alpha_0(\tau)-\widehat{\alpha}_r(\tau)\right], \\
    \operatorname{MAE}(\tau)
      &=\frac{1}{R}\sum_{r=1}^{R}\left|\widehat{\alpha}_r(\tau)-\alpha_0(\tau)\right|, \\
    \operatorname{RMSE}(\tau)
      &=\left\{\frac{1}{R}\sum_{r=1}^{R}\left[\widehat{\alpha}_r(\tau)-\alpha_0(\tau)\right]^2\right\}^{1/2}.
  \end{aligned}
  \label{eq:performance-point}
\end{equation}
Measures are computed separately by estimator and design cell. Positive Bias means estimates lie below the truth on average; MAE records absolute error, and RMSE gives greater weight to large errors.

Coverage is determined by evaluating the test statistic directly at the exact true parameter. With $c_{0.95}$ denoting the critical value defined below, define
\begin{equation}
  C_r(\tau)
    =\mathbf{1}\!\left\{W_{N,r}\!\left(\alpha_0(\tau)\right)\leq c_{0.95}\right\},
  \qquad
  \widehat{\operatorname{Coverage}}(\tau)
    =\frac{1}{R}\sum_{r=1}^{R}C_r(\tau).
  \label{eq:performance-coverage}
\end{equation}
Coverage evaluates $W_{N,r}$ directly at $\alpha_0(\tau)$ and is therefore independent of the numerical grid. Empirical size is
\begin{equation}
  \widehat{\operatorname{Size}}(\tau)
    =\frac{1}{R}\sum_{r=1}^{R}
      \mathbf{1}\!\left\{W_{N,r}\!\left(\alpha_0(\tau)\right)>c_{0.95}\right\}
    =1-\widehat{\operatorname{Coverage}}(\tau).
  \label{eq:performance-size}
\end{equation}
The equality follows from the implemented rejection rule.

Power is evaluated at deliberately false treatment-effect values
\begin{equation}
  a_{\Delta}(\tau)=\alpha_0(\tau)+\Delta,
  \qquad
  \Delta\in\{-1.00,-0.50,-0.25,0.25,0.50,1.00\}.
  \label{eq:performance-alternatives}
\end{equation}
For each offset, empirical power is
\begin{equation}
  \widehat{\operatorname{Power}}(\Delta,\tau)
    =\frac{1}{R}\sum_{r=1}^{R}
      \mathbf{1}\!\left\{W_{N,r}\!\left(a_{\Delta}(\tau)\right)>c_{0.95}\right\}.
  \label{eq:performance-power}
\end{equation}
The statistic is evaluated directly at each false value, independently of the confidence-region grid; $\Delta=0$ is size, not power.

The accepted-set measure totals all accepted components inside the primary computational domain, including disconnected components. It is neither the span between extreme accepted values nor an unrestricted confidence-region length. Because $d_z=2$, the 95\% cutoff for the statistic in Section~3.4 is
\begin{equation}
  c_{0.95}=\chi^2_{2,0.95}=5.991464547\ldots.
  \label{eq:wn-critical-value}
\end{equation}
A candidate is accepted when $W_N(a)\leq c_{0.95}$. This asymptotic cutoff is not an exact finite-sample guarantee. The primary grid is
\begin{equation}
  \begin{aligned}
    A_0&=[-1,3],
    &h&=0.1,\\
    G_0&=\{a_j=-1+(j-1)h:j=1,\ldots,41\},\\
       &=\{-1,-0.9,\ldots,2.9,3\}.
  \end{aligned}
  \label{eq:primary-grid}
\end{equation}
The same grid is used throughout. $A_0$ is a computational domain, not an asserted parameter space, so reaching its boundary establishes neither global boundedness nor unboundedness.

The initial numerical design considered reporting total confidence-region length after adaptive domain expansion. Numerical grid diagnostics showed that, under weak relevance, accepted sets could remain non-localized over increasingly wide numerical domains. The final design therefore reports a fixed-domain accepted-set measure and treats domain expansion as a separate numerical diagnostic rather than assigning an artificial finite confidence-region length.

For replication $r$, define the binary grid-acceptance indicator, accepted grid, and all-grid-points indicator by
\begin{equation}
  \begin{aligned}
    I_r(a_j)
      &=\mathbf{1}\!\left\{W_{N,r}(a_j)\leq c_{0.95}\right\}, \\
    \mathcal{C}_{r,G_0}
      &=\left\{a_j\in G_0:I_r(a_j)=1\right\}, \\
    F_r
      &=\mathbf{1}\!\left\{I_r(a_j)=1\ \text{for every }a_j\in G_0\right\}.
  \end{aligned}
  \label{eq:grid-acceptance}
\end{equation}
The frozen implementation computes the grid-based accepted-set measure within $A_0$ as the trapezoidal integral of this binary indicator:
\begin{equation}
  M_r(A_0)
  =\sum_{j=1}^{40}\frac{a_{j+1}-a_j}{2}
    \left[I_r(a_j)+I_r(a_{j+1})\right].
  \label{eq:accepted-measure}
\end{equation}
The trapezoidal sum includes separated accepted components without filling rejected gaps and does not interpolate critical-value crossings. It ranges from zero to four. The Monte Carlo mean of $F_r$ is the frequency with which all tested grid points are accepted, not acceptance of the real line.

Coverage and power remain direct evaluations outside this grid construction. A post-main diagnostic examines sensitivity to the finite endpoints using
\begin{equation}
  A_0=[-1,3],\qquad A_1=[-3,5],\qquad A_2=[-5,7].
  \label{eq:grid-diagnostic-domains}
\end{equation}
It uses $n=1000$, 100 saved primitive samples, all five quantiles and estimators, $\kappa\in\{1,0.25,0.10\}$, and spacing $0.1$. Stage~1 evaluates $A_1$ and its nested $A_0$ profile; Stage~2 reloads the same samples and adds only the tails needed for $A_2$. Domain comparisons therefore change candidates, not samples. The outputs record accepted measure, grid share and count, boundary acceptance, all-grid acceptance, and incremental measure.

Boundary acceptance shows only a lack of localization within the examined domain; it does not establish global unboundedness, acceptance of the real line, or failed identification. The bounded diagnostic ends at $A_2$, while main summaries retain $A_0$ and direct coverage and power remain unchanged.
